\chapter{导论：光学在三维重建里管什么事}
\label{ch:00}

\section{效果}
\label{sec:ch00-effect}

读完本章，你应能回答下列问题：
\begin{enumerate}
  \item 从「照明 / 投影」到「点云」的链路里，光学在哪些环节介入，哪些误差看似算法、实则来自光路与成像；
  \item 几何光学、波动光学、光度学三条线索各自管什么事，以及它们如何映射到相机内参 $K$、畸变、条纹相位 $\varphi$、信噪比 $\SNR$ 等工程量；
  \item 本书的阅读地图（\chapref{ch:00}--\chapref{ch:10}）与每章「为什么要懂」的落点；
  \item 在医学三维重建（口内扫描、结构光、近距离立体）场景中，如何用误差预算的思路做诊断，而不是只调算法超参。
\end{enumerate}

\section{原理}
\label{sec:ch00-principle}

\subsection{重建链路：光子如何变成点}

一次结构光或立体视觉重建，可以粗略写成：
\begin{center}
光源/投影 $\to$ 场景表面 $\to$ 镜头 $\to$ 传感器 $\to$ 图像/相位 $\to$ 几何解算 $\to$ 点云.
\end{center}
计算机视觉教材常从针孔相机与三角测量讲起 \citep{HartleyZisserman2004}；光学教材则从光线、波前与辐射度量讲起 \citep{HechtOptics}。本书把两边接起来：投影模型里的焦距 $f$、像面坐标 $(u,v)$、世界坐标 $(X,Y,Z)$，并不是算法里随便起的名字，它们对应可检验的光路假设。

把误差按环节拆开，比笼统说「重建不准」更有用：
\begin{itemize}
  \item \textbf{几何环节}：针孔/薄透镜是否成立、畸变是否建模、基线与景深是否匹配工作距；
  \item \textbf{波动 / 干涉环节}：条纹对比度、衍射极限、散斑与相干；相位 $\varphi$ 与光程差 $\OPD$（或 $\opd$）如何耦合；
  \item \textbf{光度 / 噪声环节}：辐照度 $E$、辐亮度 $L$、辐射通量 $\Phi$ 与相机灰度的关系；$\SNR$ 如何限制相位与匹配精度。
\end{itemize}
点云上的毫米级或微米级偏差，往往是上述环节的放大结果，而不是单一公式写错。

\subsection{三条光学线索}

\textbf{几何光学}关心「光线往哪走」。针孔给出直线投影；薄透镜公式
\[
\frac{1}{s}+\frac{1}{s'}=\frac{1}{f}
\]
把物距 $s$、像距 $s'$ 与焦距 $f$ 绑在一起。记号锁定：$f$ \textbf{只}表示焦距；光学频率用 $\nuopt$，不用 $f$ 表示频率。景深相关量记为 $Z_d$，避免与世界坐标深度 $Z$ 混用。畸变使直线变弯，必须在投影模型里显式出现（\chapref{ch:01}、\chapref{ch:02}）。

\textbf{波动光学}关心「相位与分辨率」。干涉条纹、衍射与点扩散函数 $\PSF$（记 $h$）、调制传递函数 $\MTF$，决定「你以为很锐的边」在传感器上实际有多糊 \citep{GoodmanFourierOptics,BornWolf}。结构光的相位展开与解缠，依赖条纹可辨与相位稳定 \citep{GhigliaPritt1998}（\chapref{ch:03}、\chapref{ch:06}）。

\textbf{光度学与辐射度}关心「有多少光、灰度代表什么」。灰度 $\neq$ 辐照度 $E$；曝光、增益、gamma 与噪声模型共同决定 $\SNR$。投影机偏光片、相机偏振与部分相干，影响对比度与散斑（\chapref{ch:04}、\chapref{ch:05}）。

\subsection{全书地图}

\begin{center}
\begin{tabular}{@{}clp{0.48\textwidth}@{}}
\toprule
章 & 标题要点 & 核心「为什么要懂」 \\
\midrule
00 & 导论 & 误差从哪来；地图 \\
01 & 几何光学与成像 & 针孔/薄透镜/景深/畸变如何进投影模型 \\
02 & 相机投影的光学约束 & 内参、畸变、景深与标定可观性 \\
03 & 波动光学够用本 & 干涉、衍射极限与分辨率 \\
04 & 偏振与部分相干 & 偏光片、对比度、散斑 \\
05 & 光度学与相机响应 & 曝光、噪声、灰度$\neq E$ \\
06 & 投影光学与结构光 & 条纹、相位、DMD/LCoS \\
07 & 立体与三角测量限度 & 基线、景深、遮挡、标定放大 \\
08 & 色散、色差、PSF/MTF & 精度与模糊预算 \\
09 & 共聚焦/OCT/ToF 对照 & 能力边界（选读） \\
10 & 光子到点云误差预算 & 综合检查清单 \\
\bottomrule
\end{tabular}
\end{center}

阅读建议：Round~A（本章至 \chapref{ch:02}）建立几何与相机语言；其后波动与光度为结构光与立体做准备；\chapref{ch:10} 把各章误差收成一张预算图。

\subsection{误差预算的直觉}

设某深度估计依赖视差或相位，粗记
\[
\delta Z \;\sim\; g(\text{几何标定},\,\text{光学模糊},\,\SNR,\,\text{算法}).
\]
$g$ 的具体形式随方法而变，但诊断顺序应固定：先问光路与标定是否在工作距内成立，再问条纹/纹理的 $\SNR$ 与对比度，最后才动匹配窗口或网络权重。把「光学」当成可计算的预算项，而不是背景知识，是本书的立场。


\subsection{符号与阅读约定（开篇锁定）}

后文反复出现的符号以 \texttt{STYLE.md} 的 Notation Lock 为准，此处只强调最易混的几组：
\begin{itemize}
  \item 波长 $\wavelength$、光频 $\nuopt$；焦距一律 $f$。写「绿光频率」时用 $\nuopt$，不要写 $f$。
  \item 像面 $(u,v)$、世界 $(X,Y,Z)$；景深相关量 $Z_d$。
  \item 强度语境下的 $I$ 易与辐照度混淆时，改用 $E$；辐亮度 $L$、通量 $\Phi$。
  \item 相位 $\varphi$（或 $\phi$）；光程差 $\OPD$ 或 $\opd$；点扩散 $\PSF$/$h$；$\MTF$；信噪比 $\SNR$。
  \item 相机内参 $K$，外参 $R,t$。
\end{itemize}
交叉引用用「第~\ref{sec:…}~节」或 `\chapref`；文献只引用 \texttt{refs/book.bib} 中已有 key。

\subsection{「算法误差」与「光学误差」如何分工}

在 C++/CUDA 实现里，常见惯性是：点云一差就调双边滤波、改匹配窗口、加网络 refinement。更稳的分工是：
\begin{enumerate}
  \item \textbf{可重复的系统偏差}（整弓弯曲、尺度随工作距漂）→ 先查 $K$、畸变、对焦是否锁定、窗口玻璃与温度；
  \item \textbf{随纹理/条纹质量变的随机散} → 查离焦、$Z_d$、曝光与 $\SNR$、偏振与反光；
  \item \textbf{仅在遮挡边缘或高曲率处坏} → 查 $\MTF$、部分像元混合、三角测量几何放大（\chapref{ch:07}、\chapref{ch:08}）。
\end{enumerate}
光学并不能替代数值算法，但能告诉你\textbf{算法不该被要求去拟合的物理欠缺}。把欠缺写进误差预算，比在损失里加一项盲目平滑更可控。


\subsection{与读者背景的对齐}

你可能更熟悉：损失函数、CUDA kernel、点云 ICP 与分割网络。光学公式不会取代这些工具，但会限制它们的输入是否物理上可辨。例如：
\begin{itemize}
  \item 分割很准也不能恢复因离焦丢失的条纹高频；
  \item ICP 很稳也不能拆掉错误畸变造成的全局弯曲；
  \item 提高网络容量不能单方面取消衍射极限或光子噪声。
\end{itemize}
因此本书每章的「对象」段都强制落到相机、投影、干涉、条纹或点云误差之一，避免停在定义罗列。有数竞与统计背景时，把 $Z_d$、$\SNR$、重投影误差当成可估计的量，比背诵透镜分类表更有用。

\section{限度}
\label{sec:ch00-limits}

\begin{itemize}
  \item 本书补「够用」光学，不是完整物理光学课程；电磁场边值、严格矢量衍射等从略，需要时指向 \citep{BornWolf,HechtOptics}。
  \item 「对象」段给的是诊断与设计检查清单，不是厂商操作手册或监管认证流程。
  \item 医学场景示例（口内、近距结构光）用于锚定误差来源；参数量级因设备而异，文中数量级仅供建立直觉。
  \item 机器学习重建若端到端拟合，仍受镜头畸变、景深与 $\SNR$ 约束；光学限度不会因换网络而消失。
\end{itemize}

\section{对象}
\label{sec:ch00-object}

以近距离医学结构光三维重建为例（口内扫描仪一类系统）：
\begin{enumerate}
  \item \textbf{工作距短、景深窄}：牙面起伏大时，离焦使条纹对比度下降，$\PSF$ 展宽直接抬高相位噪声；此时应查 $Z_d$ 与光圈，而不是只加大平滑。
  \item \textbf{湿润高反与次表面散射}：光度模型偏离朗伯，$E$ 与灰度非线性；偏振与波长选择影响有效对比度（后文 \chapref{ch:04}、\chapref{ch:05}）。
  \item \textbf{手持运动与标定}：外参 $R,t$ 与内参 $K$ 的微小误差经三角测量放大；先区分「标定漂移」与「算法阈值」。
  \item \textbf{检查清单起手式}：点云局部糊或系统性弯——先看畸变与 $K$；条纹区噪声大——先看离焦与 $\SNR$；深度跳变在边缘——先看 $\MTF$ 与遮挡。
\end{enumerate}


\section*{本章收束}
光学在三维重建里管的不是「背景知识」，而是误差预算里可指认的几项：几何（$f$、$K$、畸变、$Z_d$）、波动（相位、$\PSF$、$\MTF$）、光度（$E$、$L$、$\SNR$）。下一章从针孔与薄透镜把几何线索写透。

\section{习题}
\label{sec:ch00-exercises}

\begin{enumerate}
  \item \textbf{（概念）} 说明为何本书规定 $f$ 只表示焦距、光学频率必须用 $\nu$。若有人在笔记里写「绿色光 $f=5.6\times10^{14}\,\mathrm{Hz}$」，会发生什么混淆？
  \item \textbf{（概念）} 列出从投影到点云的链路中不少于四个光学相关环节，并各举一种典型误差表现（例如点云弯曲、空洞、高频抖动）。
  \item \textbf{（重建场景）} 口内扫描时咬合面中心清晰、邻面沟窝点云发散。优先怀疑景深/离焦还是匹配窗口过大？写出你的诊断顺序（不少于三步）。
  \item \textbf{（概念）} 灰度图像很亮，是否意味着辐照度 $E$ 一定高？举一个反例场景（涉及曝光或增益）。
  \item \textbf{（地图）} 若不看波动光学，直接做结构光相位，会缺哪类「限度」认知？指出本书哪一章补上。
  \item \textbf{（计算直觉）} 某系统深度误差粗估 $\delta Z \approx k/\SNR$。若 $\SNR$ 从 $20$ 升到 $40$（线性比），$k$ 不变，$\delta Z$ 如何变？这能否代替几何标定改进？为何？
\end{enumerate}

\section{习题讲解}
\label{sec:ch00-solutions}

\begin{enumerate}
  \item $f$ 在成像与相机模型中几乎总是焦距（像素焦距或毫米焦距）。用 $f$ 记频率会与薄透镜公式、内参 $f_x,f_y$ 冲突。应写 $\nu\approx5.6\times10^{14}\,\mathrm{Hz}$，波长仍用 $\wavelength$。
  \item 示例：投影光学（条纹对比度低→相位噪）；镜头畸变（直线牙弓被弯）；离焦/$\PSF$（边缘点云糊）；传感器噪声/$\SNR$（深度抖动）；标定中 $K,R,t$（全局尺度或弯曲）。任选四点并对应表现即可。
  \item 优先离焦/景深：沟窝相对对焦面更远，$Z_d$ 不够则条纹 MTF 下降。顺序建议：(1) 查对焦面与工作距是否落在声称景深内；(2) 看原图条纹对比度是否随深度变差；(3) 再查是否运动模糊或匹配参数过平滑。窗口过大通常造成过平滑而非「仅沟窝发散」。
  \item 否。提高模拟增益或延长曝光可使灰度饱和，而场景 $E$ 未变；或 tonemap/gamma 改变灰度。灰度是相机响应后的数字量。
  \item 缺少干涉/衍射/分辨率与条纹可辨的限度，对应 \chapref{ch:03}（及 \chapref{ch:06} 的相位实践）。
  \item $\SNR$ 加倍则 $\delta Z$ 减半。不能代替标定：系统偏差（错误的 $K$ 或畸变）不随 $\SNR$ 消失；$\SNR$ 主要压随机误差。
\end{enumerate}
